\chapter{Implementation and Results} 
\label{chap:implementation_results} 
% TODO-RENAME: RQ-0xx (security requirements) collides with RQ1/RQ2/RQ3 (research questions).
% Rename to SR-0xx across ALL chapters, appendices, tables and \label keys in a single coordinated pass.
This chapter applies the methodology of Chapter \ref{chap:methodology} to a physical System Under Test. It documents the implemented environment, the application of the method, the collected evidence, the verification of the derived requirements, the validation of the methodology against the thesis objective, and the reliability of the findings.

The practical work covers the direct-access variant \gls{acr:NV1} defined in Section \ref{subsec:hybrid_approach}. The passive-observation variant \gls{acr:NV2} and the in-path variant \gls{acr:NV3} are part of the methodology but are not realised in the laboratory. The methodology is therefore validated only in part. Its central steps are exercised on a real device, whereas the tests that depend on \gls{acr:NV2} or \gls{acr:NV3} remain specified but not executed.

The chapter follows the evidence-first principle of Section \ref{subsec:verdicts}. An activity is reported as performed only when a stored artefact documents it. A tool execution is not treated as requirement verification. An observation is kept separate from its interpretation. Expected behaviour from a test specification is not presented as an observed result.





\section{Scope of the Practical Demonstration}
\label{sec:testbed_setup}

\subsection{System Under Test}
\label{subsec:test_object} 

The System Under Test is the Safe Remote \gls{acr:I/O} (\gls{acr:SRIO}), a fieldbus-based safety component that exchanges safety-related input and output data with a higher-level safety controller over PROFINET and PROFIsafe. The evaluated unit identified itself over the IoT-Core interface as product \texttt{ifm-AL400S}, hardware version 1.4.0.2, hardware revision AA, and software revision 1.1.4.888. All reported results refer exclusively to this configuration.

Following the asset classification of Section \ref{subsec:assets_attack_surfaces}, the \gls{acr:SRIO} is decomposed into hardware, software, and data assets. For the practical assessment, the externally accessible attack surface was divided into the \gls{acr:HTTP}-based \gls{acr:IoT}-Core diagnostic interface, the \gls{acr:PROFINET} and PROFIsafe communication path, the \gls{acr:DCP} configuration functions, and the firmware update mechanism. Only the interfaces accessible from the implemented \gls{acr:NV1} position were exercised. Internal debug and inter-processor interfaces, such as the \gls{acr:SysCom} link between the \gls{acr:COM} and \gls{acr:SCPU} modules and the \gls{acr:JTAG} interface, are included in the asset model but were not accessible within the practical test scope.

Table \ref{tab:assets} summarises the assets and their protection objectives. It is not exhaustive, but it covers the main elements that are relevant for the security assessment.
{\renewcommand{\arraystretch}{1.2}
	\small 
	\begin{xltabular}{\textwidth}{@{} 
			c 
			>{\raggedright\arraybackslash\hsize=0.9\hsize}X 
			>{\raggedright\arraybackslash\hsize=1.4\hsize}X 
			>{\raggedright\arraybackslash\hsize=0.7\hsize}X 
			@{}}
		
		% --- Kopf- und Fußzeilen-Definition für die Longtable ---
		
		% 1. Kopfzeile auf der ersten Seite (inkl. Caption)
		\caption{Overview of assets and protection objectives}
		\label{tab:assets} \\
		\toprule
		\textbf{ID} & \textbf{Asset} & \textbf{Description} & \textbf{Protection \newline objectives} \\ 
		\midrule
		\endfirsthead
		
		% 2. Kopfzeile auf allen folgenden Seiten
		\caption*{Overview of assets and protection objectives (continued)} \\
		\toprule
		\textbf{ID} & \textbf{Asset} & \textbf{Description} & \textbf{Protection \newline objectives} \\ 
		\midrule
		\endhead
		
		% 3. Fußzeile auf allen Seiten außer der letzten
		\midrule
		\multicolumn{4}{r}{\footnotesize\textit{Continued on next page}} \\
		\endfoot
		
		% 4. Fußzeile auf der allerletzten Seite der Tabelle
		\bottomrule
		\endlastfoot
		
		% --- Eigentlicher Tabelleninhalt ---
		
		A & Trusted Safety Function & SRIO correctly executes its specified safety function. & Integrity, \newline Availability \\ 
		
		B & Integrity of Safety Configuration & Safe behaviour is determined by the intended configuration. & Accountability, \newline Authorisation, \newline Integrity, \newline Authenticity \\ 
		
		C & Integrity and Authenticity of Safety-Relevant Process Data & Input, output and PROFIsafe data correspond to the actual safety state. & Integrity, \newline Authenticity, \newline Availability \\ 
		
		D & Authenticity and Integrity of Safety Software & Bootloader and firmware remain authentic and unmodified. & Accountability, \newline Authenticity, \newline Integrity \\ 
		
		E & Integrity of Safety Monitoring & Self-tests, plausibility checks, diagnostics and fault responses remain trustworthy. & Integrity, \newline Availability \\ 
		
		F & Separation of Safety and Non-Safety Domain & A compromised COM system must not affect the integrity of the safety function. & Confidentiality, \newline Integrity, \newline Availability \\ 
		
		G & Integrity of Operating Mode & Test and update functions must not be activated or used without authorisation. & Accountability, \newline Authorisation, \newline Integrity, \newline Authenticity, \newline Availability \\ 
		
		H & SRIO Functionality & Availability of SRIO functionality shall be ensured. & Availability \\ 
		
		I & Integrity, Availability and Confidentiality of Audit/Tracing Data & Evidence of interventions, configuration changes and installed software versions must be generated, retained for the mandated period and access-restricted. & Accountability, \newline Availability, \newline Confidentiality \\
		
	\end{xltabular}
}

\subsection{Implemented Test Environment}
\label{subsec:realized_environment}
The test environment used one switch and one flat address range. Figure \ref{fig:testbed_topology} shows the physical topology, and Table \ref{tab:testbed_components} lists the components and their roles.

\begin{table}[H]
	\centering
	\small
	\caption{Components of the implemented test environment}
	\label{tab:testbed_components}
	\begin{tabularx}{\textwidth}{l l c X}
		\toprule
		\textbf{Component} & \textbf{Address} & \textbf{Port} & \textbf{Role} \\
		\midrule
		Industrial switch & 192.168.0.91 & -- & Layer-2 connectivity for all nodes \\
		PLC (F-Host) & 192.168.0.1 & 1 & Safety controller and PROFINET/PROFIsafe counterpart \\
		Safe Remote I/O (SUT) & 192.168.0.2 & 2 & Device under test \\
		Kali Linux test system & 192.168.0.80 & 3 & Security test system, connected through a USB 3.0 to Gigabit Ethernet adapter \\
		Company laptop & 192.168.0.7 & 4 & Engineering station providing TIA Portal for planned identification and maintenance read operations \\
		\bottomrule
	\end{tabularx}
\end{table}

All nodes shared the 192.168.0.0/24 segment and one Ethernet broadcast domain. No \gls{acr:VLAN} separation was configured. The test system reached the \gls{acr:SUT} and the controller directly over \gls{acr:IP}. The test object was the \gls{acr:SRIO} on port 2. The test initiator was the Kali system on port 3. The company laptop on port 4 served as the engineering station and provided \gls{acr:TIA} Portal for planned identification and maintenance read operations. These read operations were not part of the evidence-backed execution subset reported in this chapter. The available observation points were the direct request and response traffic of the test system and the \gls{acr:HTTP} responses of the \gls{acr:IoT}-Core service. Traffic that was not addressed to the test system was not necessarily visible, because a switch forwards unicast frames only to the destination port. \gls{acr:NV1} was therefore the realised position. \gls{acr:NV2} and \gls{acr:NV3} were not implemented.

The device documentation defines organisational preconditions for secure operation. The intended use assumes the device in a protected zone behind a firewall. A firmware update is expected to be performed by a trusted person behind that firewall. A security risk assessment is required. The debug interface is primarily a physical interface. These preconditions bound the threat model of the practical tests.
These preconditions are relevant to the compensating-countermeasure allocation: zone protection, trusted update procedures and the system-level security risk assessment may address obligations that the component does not implement itself. They therefore inform the interpretation of component-capability findings in Section \ref{sec:requirement_verification}.


\begin{figure}[H]
	\centering
	\begin{tikzpicture}[
		% Definition der Design-Stile für die Bachelorarbeit
		node distance=1cm and 1cm,
		box/.style={
			draw=black!70, 
			thick, 
			rectangle, 
			rounded corners=1.5mm, 
			minimum width=4.2cm, 
			minimum height=1.8cm, 
			align=center, 
			drop shadow={opacity=0.15}, 
			fill=white,
			font=\small
		},
		switch/.style={
			box, 
			fill=gray!10, 
			minimum height=2cm, 
			minimum width=4.5cm
		},
		link/.style={
			draw=black!80, 
			thick, 
			<->, % Bidirektionale Netzwerkverbindung
			>=Latex
		},
		port/.style={
			fill=white, 
			inner sep=2pt, 
			font=\scriptsize\bfseries,
			text=black!80
		}
		]
		
		% Zentrales Element: Der Switch
		\node[switch] (switch) {
			\textbf{Industrial Switch}\\
			192.168.0.91\\[1.5mm]
			\textit{\scriptsize Layer-2 connectivity}
		};
		
		% Periphere Knoten (Nodes)
		\node[box, above left=of switch] (plc) {
			\textbf{PLC (F-Host)}\\
			192.168.0.1\\[1.5mm]
			\textit{\scriptsize Safety controller}
		};
		
		\node[box, above right=of switch] (sut) {
			\textbf{Safe Remote I/O (SUT)}\\
			192.168.0.2\\[1.5mm]
			\textit{\scriptsize Device under test}
		};
		
		\node[box, below left=of switch] (kali) {
			\textbf{Kali Linux Test System}\\
			192.168.0.80\\[1.5mm]
			\textit{\scriptsize Security test system}
		};
		
		\node[box, below right=of switch] (laptop) {
			\textbf{Company Laptop}\\
			192.168.0.7\\[1.5mm]
			\textit{\scriptsize Engineering station}
		};
		
		% Verbindungen (Kabel) ziehen und Ports am Switch beschriften
		% pos=0.75 platziert das Port-Label nahe am Switch
		\draw[link] (plc.south east) -- (switch.north west) 
		node[pos=0.75, port, xshift=-0.1cm, yshift=0.2cm] {Port 1};
		
		\draw[link] (sut.south west) -- (switch.north east) 
		node[pos=0.75, port, xshift=0.1cm, yshift=0.2cm] {Port 2};
		
		\draw[link] (kali.north east) -- (switch.south west) 
		node[pos=0.75, port, xshift=-0.1cm, yshift=-0.2cm] {Port 3};
		
		\draw[link] (laptop.north west) -- (switch.south east) 
		node[pos=0.75, port, xshift=0.1cm, yshift=-0.2cm] {Port 4};
		
	\end{tikzpicture}
	\caption{Logical topology and addressing of the implemented test environment}
	\label{fig:testbed_topology}
\end{figure}

\subsection{Practical Coverage and Limitations}
\label{sec:practical_scope}

The demonstration applied the methodology to the Safe Remote \gls{acr:I/O} under the \gls{acr:NV1} network position defined in Section \ref{subsec:hybrid_approach}. The test system acted as a direct \gls{acr:IP} participant on the shared switch. The passive-observation variant \gls{acr:NV2} and the in-path variant \gls{acr:NV3} were not realised. No mirror or \gls{acr:SPAN} port was configured, because the available laboratory switch does not support port mirroring, and the test system was never placed inline, because rebuilding the setup for an inline bridge exceeded the available time. The cyclic \gls{acr:PROFINET} and PROFIsafe communication between the controller and the \gls{acr:SUT} was therefore not observed, and no claim is made about that traffic.

The empirical coverage was bounded by the single test object, the implemented topology, the available interfaces, and the \gls{acr:NV1}-compatible activities. Within \gls{acr:NV1}, the evaluation focused on the \gls{acr:IoT}-Core \gls{acr:HTTP} interface and on tests in which the test system acted as an unauthorised communication partner. Where a company tool assumed an interface that the device does not expose, the test was recorded as an applicability finding rather than as a failure, following Section \ref{subsec:company_framework}. The scope of this practical validation must be distinguished from the achievement of the thesis objective, which is to develop and validate a systematic methodology for verifying security functions rather than to execute every technically possible test.




\section{Instantiation of the Methodology}
\label{sec:instantiation}
The applicability assessment defined in Section \ref{subsec:company_framework} was applied before each test was selected. The assessment considered the target interface, the required network position, the tool compatibility, and the available observation method. Each requirement was linked to a component requirement, a target asset, and one or more test cases. Each practical activity is described once in Section \ref{sec:executed_tests} and is referenced from both the company identifier and the thesis-specific identifier.

\subsection{Requirement, Asset and Interface Mapping}
\label{subsec:req_asset_interface}
The instantiation binds the abstract mapping of Section \ref{sec:asset_threat_function_mapping} to the concrete SRIO. Each requirement is paired with the affected asset, the security function that treats it, the interface on which it is exercised, and the network access variant that the corresponding test requires. Table \ref{tab:req_asset_interface} shows representative rows across the range of access requirements; the full requirement derivation is documented in Appendix \ref{app:requirement_derivation}.

\begin{table}[H]
	\centering
	\small
	\caption{Representative instantiation of the requirement-to-test mapping for the SRIO}
	\label{tab:req_asset_interface}
	\begin{tabularx}{\textwidth}{l X X X l}
		\toprule
		\textbf{Req.} & \textbf{Asset} & \textbf{Security function} & \textbf{Interface} & \textbf{Access} \\
		\midrule
		RQ-003 & Audit data & Logging and auditing & IoT-Core error log & NV1 \\ \addlinespace
		RQ-005 & Safety firmware & Integrity protection & Firmware update service & NV1 \\ \addlinespace
		RQ-010 & Device availability & Availability and resilience & HTTP and safety communication & NV1 and NV2 \\ \addlinespace
		RQ-011 & Safety configuration & Authorisation & Configuration interfaces & NV1 \\
		\bottomrule
	\end{tabularx}
\end{table}

\subsection{Applicability Assessment}
\label{subsec:cm_applicability}
Table \ref{tab:cm_applicability} records the applicability outcome for the ten company tests. Not-applicable outcomes are properties known in advance from the device or topology, not test failures. Test CM-002 requires a second network zone, which the flat segment does not provide, and CM-006 requires an operating-system shell, which the SUT does not expose. CM-003 and CM-005 were executed to establish the absence of an authentication endpoint and a TLS endpoint; their negative results are therefore recorded as evidence-based applicability outcomes.

\begin{table}[htbp]
	\centering
	\small
	\caption{Applicability of the company framework tests to the SUT}
	\label{tab:cm_applicability}
	\begin{tabularx}{\textwidth}{l l X}
		\toprule
		\textbf{CM} & \textbf{Tool} & \textbf{Applicability outcome} \\
		\midrule
		CM-001 & Nmap & Applicable; performed \\
		CM-002 & Nmap & Not applicable to the topology; no zone or conduit model in a flat segment \\
		CM-003 & Hydra & Executed with negative result; no authentication endpoint, absence demonstrated by tool output \\
		CM-004 & OWASP ZAP & Applicable with adaptation; web scan performed, multi-role comparison not applicable \\
		CM-005 & SSLscan & Executed with negative result; no TLS endpoint, absence demonstrated by tool output \\
		CM-006 & Lynis & Not applicable to the device; no operating-system shell \\
		CM-007 & Nikto & Applicable; performed \\
		CM-008 & OpenVAS/GVM & Blocked; scanner could not be made operational in the laboratory setup \\
		CM-009 & Siege & Applicable; performed, load summary recorded; 50-connection run outside the specified envelope \\
		CM-010 & boofuzz & Applicable; performed \\
		\bottomrule
	\end{tabularx}
\end{table}

\subsection{Selected Test Cases and Execution Status}
\label{subsec:cm_tc_map}
Table \ref{tab:cm_tc_map} maps the company tests to the derived requirements and to the security category in which they are reported. One practical activity can support several requirements. The mapping is not one-to-one. For example, one Nmap run supports the exposure countermeasure CM-001 and the identification and asset-discovery category, and one query of the identification interface supports several identification requirements. The tests CM-002, CM-006, and CM-008 do not appear in the mapping, because they are not applicable or could not be operated on this device and therefore support no requirement. 

The table also lists the requirements for which no company test was executed under NV1, marked with a dash and a not-executed evidence entry. RQ-002 and RQ-005 belong to the firmware and configuration-integrity category (Section \ref{subsec:cat_config_firmware}), and RQ-009 requires an induced configuration change that was not performed (Section \ref{subsec:cat_logging}).
\begin{table}[H]
	\centering
	\small
	\caption{Mapping of company tests, requirements, and security categories}
	\label{tab:cm_tc_map}
	\begin{tabularx}{\textwidth}{l l X X}
		\toprule
		\textbf{CM} & \textbf{Req.} & \textbf{Test category} & \textbf{Evidence} \\
		\midrule
		CM-001, CM-007 & RQ-001 & Identification and asset discovery & Nmap, Nikto \\
		- & RQ-004, RQ-006, RQ-007, RQ-013 & Identification and asset discovery & IoT-Core inventory \\
		CM-005 & RQ-004, RQ-014 & Communication security & SSLscan \\
		- & RQ-003, RQ-008, RQ-012 & Logging and traceability & error log, systick \\
		CM-009 & RQ-007, RQ-010 & Availability and resilience & Siege \\
		CM-010 & RQ-010 & Robustness & boofuzz \\
		CM-003, CM-004 & RQ-011, RQ-014 & Authorisation and access control & Hydra, ZAP, IoT-Core \\
		- & RQ-002, RQ-005 & Firmware and configuration integrity & none not executed \\
		- & RQ-009 & Logging and traceability & not executed \\
		\bottomrule
	\end{tabularx}
\end{table}

\section{Test Execution and Results}
\label{sec:executed_tests}

The two-level test catalogue derived in Chapter \ref{chap:methodology} defines abstract test scenarios for all fourteen requirements, including the tests that depend on the passive-observation variant NV2 and the in-path variant NV3. Only a subset of these was executed in the practical phase: the applicable company tests and the requirement-oriented tests that the NV1 position supports. This section groups the tests into security categories and describes them at the level of the test activity, the method, and the result, rather than by individual test-case identifier. Each category summarises the abstract scenarios it contains, and then reports the concrete execution and the stored evidence for the activities performed under NV1. The tests that require NV2 or NV3 are described as designed but not executed. Every reported result is limited to what the stored evidence shows. A tool execution is not treated as requirement verification, and an expected outcome from a test specification is not presented as an observed result.

Each category follows the same logic. It states the machine-level obligation from Regulation (EU) 2023/1230 from which the component requirement is derived, what must therefore be checked, how the property can be tested and which company countermeasure applies, and what the general finding on the device was. Every performed test is stored in some form as evidence, depending on the test as a text file, an HTML report, a log file, or a screenshot. The exact tool commands for the executed activities are listed in Appendix \ref{app:tool_commands}.

\subsection{Identification and asset discovery}
\label{subsec:cat_identification}

This category covers the discovery of the reachable attack surface and the identification of the installed software and data, addressing RQ-001 on the non-safety path and requirements RQ-004, RQ-006, RQ-007, and RQ-013. It targets the separation of the safety and non-safety domains and the software and data-identification assets. The abstract catalogue contains the attack-surface enumeration, the identification-surface inventory, the software-bill-of-materials completeness analysis, the plaintext-transport observation, the data-signature uniqueness check, and the version-history and per-upload-timestamp checks. The bill-of-materials, signature-uniqueness, and version-history scenarios were designed but not executed. The activities executed under NV1 are reported below.

\paragraph{Attack-Surface Enumeration}
\label{subsec:res_attack_surface}
This activity examined the externally reachable services on the non-safety path. It relates to CM-001 and CM-007, to requirement RQ-001, and to asset F, the separation of the safety and the non-safety domain. It targets the IoT-Core HTTP path and was run under NV1. The commands are listed in Appendix \ref{app:cmd_attacksurface}.

A full TCP port scan and a top-ports UDP scan were run against the SUT with Nmap. The TCP scan covered all 65535 ports and reported one open port, port 80, which returned HTTP responses without a server banner. All other TCP ports were reported as closed. The open HTTP port is consistent with the IoT-Core diagnostic interface. The UDP scan of the top 200 ports reported port 161 as open, identified as SNMP, and port 49152 as open or filtered. An open SNMP port is consistent with a PROFINET device, because the PROFINET standard uses SNMP for network diagnostics. The reported open-port count depends on the scan timing. At the slower \texttt{-T3} timing the scan reported the stable result above, whereas the faster \texttt{-T5} timing reported further open ports. The additional ports did not reappear on repetition, which suggests that the faster timing affects the observed responses rather than revealing additional stable services. Packet loss, SYN-backlog exhaustion, rate limiting, device overload, or test-system effects are alternative explanations. If aggressive timing does destabilise the device's responses, that would itself be an availability observation relevant to RQ-010; this remains an open item.

The externally reachable services on the non-safety path are the HTTP service on port 80 and the SNMP service on port 161. Nikto targets the exposure of debug and service interfaces on a web server, which is the objective of the company countermeasure CM-007. On this device the physical debug interface is encapsulated and not reachable, so the scan characterises only the reachable HTTP service surface. The Nikto report lists missing HTTP security headers, including content-security-policy, x-content-type-options, strict-transport-security, and permissions-policy. These headers are a defence-in-depth measure, so their absence is an observation rather than a direct vulnerability. The report also flags several \texttt{/actuator/*} paths as Spring Boot Actuator endpoints that return JSON.

The Nikto findings are unvalidated scanner results \cite[p. 39]{richier2011security}. The Spring Boot Actuator finding is a probable false positive. The device is an ifm IoT-Core and not a Spring Boot application, and its HTTP service returns a generic JSON structure for many paths, which a generic scanner signature can misclassify as an exposed Actuator endpoint. The finding therefore does not, on its own, indicate a risk. Directly requesting one such path and inspecting the raw response remains an open verification item; no result is claimed here.

The enumeration supports RQ-001 at the level of the non-safety path. It shows which services a connecting device can reach. It does not, on its own, demonstrate a hazardous situation. The SNMP service on port 161 was deliberately excluded from the practical scope because it was not examined beyond enumeration. Its authentication, information-disclosure and logging implications remain open items for a future test campaign.


\paragraph{Identification-Surface Inventory}
\label{subsec:res_identification}
This activity enumerated the identification data that the device exposes. It relates to requirements RQ-004, RQ-006, RQ-007, and RQ-013, and to assets B and D. It targets the IoT-Core \texttt{deviceinfo} and \texttt{firmware} paths and was run under NV1. The commands are listed in Appendix \ref{app:cmd_inventory}.

The identification interface was queried over HTTP. The interface required no credentials, so the read operation was anonymous. For each queried path the device returned a JSON response with status code 200 that contained the requested identifier. The read operations covered the hardware version and revision, the software revision, the COM and SCPU firmware and bootloader versions, the host firmware version, and the fieldbus firmware version. A successful read operation is recognisable by the returned status code 200, which indicates a successful response. A separate single read of the software revision also returned status 200, which documents the baseline availability of the interface.

The device exposes top-level version identifiers for the COM firmware, the SCPU firmware, the bootloaders, and the fieldbus firmware, all through anonymous read access. The inventory supports RQ-004, RQ-006, and RQ-007 at product granularity. It confirms that safety-relevant software is identifiable and available on request. It does not establish a complete software bill of materials (SBOM). The versions of third-party components, such as the communication stack or the real-time operating system, are not individually exposed. The visible data is therefore an identification surface rather than a complete installed-software inventory or a full SBOM.


\subsection{Communication security}
\label{subsec:cat_communication}
This category covers the protection of the communication paths against corruption and disclosure, addressing RQ-001, RQ-002, and the transport aspect of RQ-014. It targets the safety-relevant process data and the audit data. The abstract catalogue contains the transport-confidentiality observation on the non-safety path, the integrity-rejection and in-path manipulation scenarios for the safety telegram, and the watchdog-passivation scenario. The safety-telegram integrity, in-path manipulation, and watchdog scenarios require the NV2 or NV3 position or the engineering tool and were designed but not executed. The transport observation executed under NV1 is reported below.

\paragraph{Cleartext Transmission of the Identification Data}
\label{subsec:res_cleartext}
This activity examined the transport protection of the identification data. It relates to CM-005, to requirements RQ-004 and RQ-014, and to asset I. It targets the IoT-Core HTTP service and was run under NV1. The commands are listed in Appendix \ref{app:cmd_cleartext}.

SSLscan verifies the TLS configuration of a TLS endpoint. It probes for a TLS handshake on a given host and port and reports the enabled protocol versions and the certificate. The tool was directed at the HTTP service on port 80. The saved output shows that no SSL or TLS protocol version is enabled and that no certificate can be parsed. The output therefore confirms that port 80 does not provide a TLS endpoint, so the service is plain HTTP without transport encryption. The identification data of Section \ref{subsec:res_identification} observed in the test system's own requests and responses was therefore transmitted in plaintext.
The absence of TLS is a device characteristic, not a failed TLS test. It is relevant to RQ-014 as a defence-in-depth observation. A confidentiality impact requires two further conditions. The transmitted data must be sensitive, and an unauthorised party must be able to observe it. The second condition depends on the NV2 position, which was not realised. The confidentiality claim is therefore limited to the cleartext transmission that the test system observed in its own requests and responses.

\subsection{Logging and traceability}
\label{subsec:cat_logging}
This category covers the evidence, accountability, and retention requirements RQ-003, RQ-008, RQ-009, and RQ-012, with the audit-data asset. The abstract catalogue contains the intervention-evidence baseline, the time-attribution and cold-start-deletion scenarios, the software- and configuration-modification evidence scenarios, and the circular-buffer and retention scenarios. The software-installation and configuration-modification evidence scenarios require an induced software or configuration change and were not fully executed. The intervention-evidence and time-attribution activity executed under NV1 is reported below; the log-access probe is reported under the authorisation and access-control category.
\paragraph{Intervention Evidence and Time Attribution}
\label{subsec:res_evidence}
This activity examined what evidence the device records for an intervention and how it is time-referenced. It relates to requirements RQ-003, RQ-008, and RQ-012, and to assets I and G. It targets the IoT-Core \texttt{errorlog} and \texttt{systick} paths and was run under NV1. The commands are listed in Appendix \ref{app:cmd_evidence}.

The machine-level requirement requires a tracing-log capability, and usable intervention evidence requires a time reference that remains unambiguous across the relevant operating states. No dedicated tracing log was exposed through the examined interface. The available substitute artefact is an error log with a relative system tick, and it is cleared at a cold restart. The activity therefore records both the apparent absence of the required tracing-log function and the properties of the substitute artefact.

The error log and the system tick counter were read over HTTP in a controlled sequence. The sequence recorded a baseline, an induced input fault by short-circuiting pins 1 and 3 of a digital input of the SRIO, a cold restart of the device by removing power, and a second induced input fault. Table \ref{tab:rq003_evidence} lists the recorded system tick values and the number of log entries at each step.

\begin{table}[H]
	\centering
	\small
	\caption{Recorded system tick values and error-log entry counts}
	\label{tab:rq003_evidence}
	\begin{tabular}{l r r}
		\toprule
		\textbf{Step} & \textbf{System tick} & \textbf{Log entries} \\
		\midrule
		Baseline & 18329 & 0 \\
		After event 1 (input fault) & 44979 & 2 \\
		After cold restart & 27509 & 0 \\
		After event 2 (input fault) & 49889 & 1 \\
		\bottomrule
	\end{tabular}
\end{table}


At the baseline the error log was empty. The short-circuit at the digital input triggered a fault, which produced two error entries on the safety CPU. Each entry carried a relative system tick value but no absolute timestamp and no actor identity. After the cold restart the error log was empty again, and the system tick had dropped, as Table \ref{tab:rq003_evidence} shows. The second input fault produced a new entry. Because the system tick is reset by a restart, an entry recorded after the restart can carry a lower tick than an earlier entry, although it occurred later in real time.

The absence of a dedicated tracing log is the primary finding for the required function. The substitute error log does collect evidence of connection-loss events, but it carries only a relative system tick and no actor identity. The system tick resets at a cold restart, so entries cannot be ordered across a power cycle. An entry recorded after the restart can carry a lower tick than an earlier entry, which confirms the reset. The log content also did not persist across the cold restart: the two entries from the first event were absent afterwards. This evidence-based observation directly demonstrates cold-start deletion in the substitute artefact and supports the conclusion that the retention capability is not demonstrated for RQ-012.


\subsection{Availability and robustness}
\label{subsec:cat_availability}
This category covers the availability aspect of RQ-007 and the resilience aspect of RQ-010, with the availability and domain-separation assets. The abstract catalogue contains the baseline resilience scenario, the network and layer-two flood scenario, the application-layer flood and domain-separation scenario, the connection-hold self-denial-of-service scenario, and the aggregate outcome-verification scenario. The application-layer load test, the connection-hold procedure, and the layer-two flood were executed under NV1. The domain-separation confirmation depends on observation of the safety channel and was designed but not executed, because it requires the NV2 position. The activities executed under NV1 are reported below; the protocol-fuzzing activity is reported under the robustness category.

\paragraph{Availability under Connection Load}
\label{subsec:res_availability}
This activity examined the availability of the identification interface under connection load. It relates to CM-009, to requirements RQ-007 and RQ-010, and to assets H and F. It targets the IoT-Core HTTP service and was run under NV1. The commands are listed in Appendix \ref{app:cmd_availability}.

A load test with Siege was run against the IoT-Core HTTP service at 50 and at 3 concurrent connections. The two load levels reflect the framework default and the device-specific expectation for the SRIO, whose requirement documentation permits at most two concurrent connections. The console output and the Siege summary logs are stored as evidence.

The acceptance criterion is bounded by the specified operating envelope: availability is assessed only for the number of concurrent connections that the product specification permits. The 3-connection run is therefore an in-envelope measurement. The 50-connection run is an overload observation outside that envelope and cannot, by itself, support a PASS or FAIL verdict.

The stored summaries show that the service responded to a large number of requests with HTTP status 200 and that a substantial number of connections failed in the same runs. Under the 50-connection load, individual runs recorded between 53 and 361 successful transactions and between 65 and 100 failed connections, and in one run no transaction completed over an interval of about 27 seconds. % TODO-NUMBER: insert exact request and failure totals, and the exact interval, from stored evidence
These results document behaviour under overload, but do not establish degraded availability within the specified operating envelope. The 3-connection run remains the relevant in-envelope evidence; its exact result is limited to the stored Siege summary.

A connection-hold procedure was also executed against the identification interface. The procedure was intended to occupy the available connection slots of the HTTP server and then issue a further request. The stored evidence shows that this additional request returned HTTP status 200, but does not establish that the connection slots were actually occupied. The negative result is therefore inconclusive and does not demonstrate resilience against self-denial of service.

The activity provides in-envelope evidence only where the 3-connection result is interpreted against the specified operating envelope. The out-of-envelope 50-connection observation and the unreproduced connection-hold result do not, by themselves, support a PASS or FAIL verdict. A statement about domain separation or about the safety function is not possible under NV1, because the safety channel was not observed.

\paragraph{Layer-Two and Network Flooding}
\label{subsec:res_l2flood}
This activity examined the resilience of the device against a link-layer and network flood. It relates to requirement RQ-010 and to assets H and F. It targets the shared Ethernet segment and was run under NV1. The commands are listed in Appendix \ref{app:cmd_flood}.

A packet flood and a MAC-address flood were initiated from the test system with hping3 and macof. During the activity, a fault indication was observed on the device status LEDs. Because no complete terminal capture and no independent observation of the safety channel were available, the activity is treated as a sample observation rather than as a conclusive resilience test. The effect on the cyclic safety channel could not be observed under NV1, because that traffic is not visible to the test system.

The device therefore entered or indicated a fault state during the disturbance on the non-safety path. The flooding tools also target the switch and the shared segment, not only the device, so the indication cannot be attributed to a device-side security response. This does not, on its own, demonstrate that a specific security function detected the flood, nor the behaviour of the safety function, because the safety channel was not monitored. To observe the effect of the flood on the cyclic safety communication, the passive-observation variant NV2 would be required. Requirement RQ-010 therefore remains inconclusive for the hazardous-situation clause at the link layer.

The robustness aspect of RQ-010 under malformed input is examined with the same availability asset. The abstract catalogue contains the protocol-fuzzing scenario for the fieldbus-discovery, safety-frame, and web-interface inputs. The fieldbus-discovery and safety-frame fuzzing require the NV2 or NV3 position and were designed but not executed. The web-interface fuzzing executed under NV1 is reported below.

\paragraph{Robustness under Malformed Input}
\label{subsec:res_robustness}
This activity examined the robustness of the HTTP service against malformed input. It relates to CM-010, to requirement RQ-010, and to asset H. It targets the IoT-Core HTTP service and was run under NV1. The commands are listed in Appendix \ref{app:cmd_robustness}.

A fuzzing test was run with boofuzz against the HTTP service on port 80. The saved script fuzzes the request line of an HTTP GET request. The boofuzz session databases and the run log are stored as evidence. The log records that most malformed requests were processed without an observable effect, but that after one malformed request the HTTP service became unreachable. % TODO-NUMBER: insert the exact number of malformed requests processed without an observable effect from stored evidence
Boofuzz could not re-open the connection for about forty seconds, and its monitor attributed the failure to the preceding test case, after which the service recovered. % TODO-NUMBER: insert the exact interruption duration from stored evidence
During the run the device signalled a bus fault and a system fault on its status LEDs, and an error was observed on the device. 

The stored evidence therefore shows a temporary disruption of the HTTP service on the non-safety path under malformed input, followed by an automatic recovery. The disruption is a finding that requires follow-up by development. The causation to a specific input, the reproducibility, and any effect on the safety function are not established, because the safety channel was not monitored. Requirement RQ-010 therefore remains inconclusive for the hazardous-situation clause, while the robustness weakness of the HTTP interface is now evidenced. The result demonstrates the value of robustness testing, because a structured fuzzing run revealed a device-side disruption that would otherwise remain undetected.

\subsection{Authorisation and access control}
\label{subsec:cat_access}
This category covers the modification-prevention and access-restriction requirements RQ-011 and RQ-014, together with the configuration-modification evidence of RQ-009, with the operating-mode, configuration, and audit-data assets. The abstract catalogue contains the read-only enforcement scenario, the unauthenticated configuration-protocol reset, the fieldbus parameter-write scenario, the firmware-update mode gate, the physical-barrier scenario, the unauthenticated tracing-log read, and the access-control and deletion-control probes. The device exposes a single anonymous read-only role and no authentication endpoint, so the authentication-oriented scenarios reduce to the observation that no identity model exists. The configuration-protocol and physical-barrier scenarios were designed but not executed. The read-only write probe and the access-control probe executed under NV1 are reported below.
\paragraph{Access Control and Configuration Interfaces}
\label{subsec:res_access_control}
This activity examined the access-control model of the non-safety interface. It relates to CM-003 and CM-004, to requirements RQ-011, RQ-014, and RQ-009, and to assets B and G. It targets the IoT-Core interface and the configuration protocol and was run under NV1. The commands are listed in Appendix \ref{app:cmd_access}.

The IoT-Core interface provided anonymous read access without a login. The Hydra procedure of CM-003 was executed with a negative result because no authentication endpoint was observed. An active web scan was run with OWASP ZAP through a local proxy, and the report was saved. A comparison of multiple user roles was not applicable, because the interface exposes a single anonymous read-only role.

A write probe was run against three IoT-Core nodes. A POST request to the safety-address node, to a digital-output node, and to the application-tag node was rejected in each case, and the responses are stored as evidence. The device answered with a bad-request response rather than an explicit authorisation denial. A bad-request response shows that the write was not accepted, but it does not by itself prove an access-control decision, because it can also result from the request format. For the application-tag node, which the device documentation lists as writable, an explicit authorisation denial would have been the expected response.

The anonymous read-only access is an observation about the IoT-Core interface. It does not, on its own, decide RQ-011 or RQ-014. The three submitted write requests did not result in an accepted modification, but the returned bad-request responses do not distinguish input validation from an authorisation decision. The activity therefore does not demonstrate an access-control mechanism and remains inconclusive with respect to RQ-011. It does not cover the full write surface, the configuration-protocol behaviour, or the factory-reset path, which were not exercised. The tracing-log access control of RQ-014 is not evidenced beyond the single anonymous read-only role. These requirements therefore remain inconclusive or only partially assessed.

\subsection{Open Test Areas}
\label{subsec:cat_config_firmware}
This subsection collects the test areas that were specified in the methodology but not executed in the practical phase, so that they are not misread as results. They comprise the firmware and configuration-integrity tests, the tests that require the NV2 or NV3 position, and the company tests that were not applicable or blocked.

The firmware and configuration-integrity tests cover the integrity and authenticity of safety-relevant software and data under RQ-005, the parameter-integrity aspect of RQ-002, and the physical tamper resistance of the address and mode interface. They target the configuration, firmware, and operating-mode assets. The abstract catalogue contains the parameter-integrity rejection and forgery scenarios, the unsigned and modified firmware installation scenario, the update-package static-analysis scenario, the analytical safety-self-test scenario, and the physical tamper-resistance scenario. These scenarios require the engineering tool, a spare unit for a risk-bearing firmware test, the NV3 in-path position, or physical access to the sealed interface. They were designed but not executed in the practical phase, so no requirement in this category is assessed from stored device evidence. The parameter-integrity, firmware-authenticity, and physical-tamper gaps identified during the theoretical validation are documented as open items for a future test campaign (Section \ref{sec:future_work}).

\label{subsec:res_na_blocked}
Three company tests could not be applied to the device, or were blocked by the environment. CM-002 requires a second zone and a conduit model, which the flat laboratory segment does not provide. CM-006 requires an operating-system shell, which the embedded device does not expose. CM-008 could not be operated, because the OpenVAS/GVM scanner could not be made operational within the laboratory setup and the available time. These outcomes are recorded as applicability or environment findings. They are meaningful methodological results, because they show why an applicability assessment is required before execution. They do not indicate a product defect, and they do not support a positive or a negative requirement verdict.

The following evidence items remain open: confirmation of the probable Spring Boot Actuator false positive by a direct request; examination of the SNMP service on port 161; verification of whether aggressive Nmap timing destabilises device responses; and repetition of the connection-hold procedure with evidence that the connection slots were occupied. These items are not treated as observed results in this chapter.
% TODO-OTHERFILE: add the unreproduced connection-hold procedure to the future work section of the summary chapter.
% TODO-OTHERFILE: add the Actuator-path confirmation to the future work section.

\section{Requirement Assessment}
\label{sec:requirement_verification}
The purpose of this section is to show that the methodology can assign a traceable verdict to each derived component-level requirement from the collected evidence. The per-requirement results are the output of applying the method to one device under the NV1 position. They are not a conformance audit of the product, and no legal conformity is claimed. A FAIL verdict denotes an evidenced absence of component capability against a derived requirement, not a declaration of legal non-conformity of the product. Under the allocation defined in \autoref{sec:regulatory_applicability}, the corresponding machine-level obligation may instead be satisfied by a compensating countermeasure at system or zone level, which the machine manufacturer or integrator must document.

This section relates the executed activities to the derived requirements. It uses the verdict scale of Section \ref{subsec:verdicts}. All activities were run under the network variant NV1. A verdict is assigned only where a stored artefact supports it. The execution status of a test and its verdict are distinct: a test can be executed and evidenced yet still leave the requirement inconclusive, and a definitive FAIL can result from an evidenced device deviation even though the test itself succeeded. Where the evidence documents such a deviation, the requirement is assessed as FAIL because the component capability is not demonstrated against the derived requirement.

Table \ref{tab:req_verification} summarises the current status per requirement. No requirement is verified as fully met. No legal conformity is claimed. A PARTIAL or an inconclusive assessment reflects the bounded scope of the practical phase.

\begin{xltabular}{\textwidth}{l X l X}
	\caption{Current requirement verification status under NV1}
	\label{tab:req_verification} \\
	\toprule
	\textbf{Req.} & \textbf{Supporting activity} & \textbf{Assessment} & \textbf{Limitation} \\
	\midrule
	\endfirsthead
	
	\caption[]{Current requirement verification status under NV1 (continued)} \\
	\toprule
	\textbf{Req.} & \textbf{Supporting activity} & \textbf{Assessment} & \textbf{Limitation} \\
	\midrule
	\endhead
	
	\midrule
	\multicolumn{4}{r}{\textit{Continued on next page}} \\
	\endfoot
	
	\bottomrule
	\endlastfoot
	
	RQ-001 & attack-surface enumeration & PARTIAL & non-safety path only; scanner findings unvalidated \\
	RQ-002 & none executed & NOT ASSESSED & requires NV2, NV3, or the engineering tool \\
	RQ-003 & intervention evidence and time attribution & FAIL & no dedicated tracing log; no absolute timestamp or actor identity; compensating control at system level required \\
	RQ-004 & identification inventory; cleartext transport & PARTIAL & no complete software bill of materials \\
	RQ-005 & none executed & NOT ASSESSED & firmware and integrity tests not executed \\
	RQ-006 & installed-software inventory & PARTIAL & product granularity only \\
	RQ-007 & identification availability; load test & PARTIAL & availability on request; 50-connection overload observation outside the specified envelope \\
	RQ-008 & software-event evidence baseline & INCONCLUSIVE & no software-installation event induced \\
	RQ-009 & none executed & NOT ASSESSED & no configuration change induced \\
	RQ-010 & HTTP fuzzing; load test; layer-two flood & INCONCLUSIVE & hazardous-situation clause needs the safety channel; flood effect not attributable to device-side response \\
	RQ-011 & write probe on three nodes & INCONCLUSIVE & bad-request responses do not prove an access-control decision \\
	RQ-012 & cold-start deletion & FAIL & substitute error log cleared by a cold restart; compensating control at system level required \\
	RQ-013 & none executed & NOT ASSESSED & version history not exposed; current-version read does not evidence retention \\
	RQ-014 & cleartext transport; access-control probe & PARTIAL & single anonymous role; exposure needs NV2 \\
\end{xltabular}

The strongest evidence concerns the intervention log. The device records connection-loss events, but no dedicated tracing log was exposed through the examined interface; the substitute error-log entries carry only a relative system tick, no actor identity, and the log content is cleared by a cold restart. The evidence therefore supports the conclusion that the required tracing capability is not demonstrated and that the substitute artefact has the documented limitations behind RQ-003 and RQ-012. The identification tests support RQ-004, RQ-006, and RQ-007 at product granularity, but not at component granularity, because no software bill of materials is exposed. RQ-013 is not assessed because a current-version read does not evidence retention of uploaded software versions. The robustness test and the load test leave RQ-010 inconclusive for the hazardous-situation clause, because the safety channel was not observed. The write probe shows only that three specific write requests were not accepted; because the device answered with a bad-request response, it does not establish an access-control decision, so RQ-011 remains inconclusive. The remaining requirements are not assessed, because the corresponding configuration, firmware, and safety-channel evidence is not available under NV1.


\paragraph{Coverage of the Practical Phase}
\label{subsec:metrics_application}
The assessment model of Section \ref{sec:evaluation_metrics} is qualitative, so the results are summarised as direct counts rather than as ratios. Fourteen requirements were derived from the regulation. Ten company tests were assessed for applicability: five were applicable, two were not applicable to the device or topology, two were executed with a negative result, and one was blocked in the laboratory. Fifteen test activities were at least partially executed under the NV1 position. Two requirements, RQ-003 and RQ-012, were assessed as FAIL on the basis of evidenced component-capability deviations. Four requirements, RQ-002, RQ-005, RQ-009 and RQ-013, were NOT ASSESSED. The remaining executed activities produced PARTIAL or inconclusive outcomes, and several requirements were not assessed because the corresponding configuration, firmware, version-history or safety-channel evidence was not collected.
Almost every executed activity is supported by stored raw evidence, including the Nmap and Nikto output, the SSLscan output, the inventory read operations, the error-log and system-tick read operations, the Siege summaries, the connection-hold output, the write-probe responses, the boofuzz session databases, and the OWASP ZAP report. The evidence base is therefore largely raw and reproducible. The small number of definitive verdicts reflects the deliberately bounded scope of the practical phase rather than a deficiency of the method. The assessment becomes fuller once the outstanding NV2, NV3, firmware, version-history and TIA-based activities are executed.

\section{Evaluation of the Methodology}
\label{sec:objective_validation}

\subsection{Validation against the Thesis Objective}
\label{subsec:objective_validation_sub}
This section validates the developed methodology against the thesis objective. The evaluation follows the evaluation activity of the design science research process \cite[p. 46]{Peffers}. The thesis objective is to develop and validate a systematic methodology that supports the verification of security functions for industrial gateways at the physical IT/OT interface.

The methodology was applied end to end on a real device. It derived requirements from the regulation, connected them to assets and interfaces, identified applicable and non-applicable tests before execution, adapted the generic company procedures, collected evidence, and related observations to technical requirements. Table \ref{tab:goal_validation} maps each objective component to the methodology element that realises it, to the practical evidence, and to the validation boundary.

\begin{table}[H]
	\centering
	\small
	\caption{Validation of the methodology against the thesis objective}
	\label{tab:goal_validation}
	\begin{tabularx}{\textwidth}{X X l}
		\toprule
		\textbf{Objective component} & \textbf{Methodology element and evidence} & \textbf{Status} \\
		\midrule
		Derive requirements from the regulation & Machinery Regulation to prEN 50742 to IEC 62443-4-2 derivation; 14 requirements (Appendix \ref{app:requirement_derivation}) & demonstrated \\
		Connect requirements to assets and interfaces & asset mapping and asset table; requirement-to-asset links & demonstrated \\
		Assess security testing approaches & structured comparison and selected combination & demonstrated \\
		Evaluate the company framework & applicability assessment; Table \ref{tab:cm_applicability} & demonstrated \\
		Identify tool-to-target mismatches & CM-003 and CM-005 executed with negative result; CM-006 not applicable; CM-008 blocked & demonstrated \\
		Define requirement-oriented test cases & two-level catalogue; executed subset & demonstrated in part \\
		Define test-environment variants & NV1, NV2, NV3 model; NV1 realised & demonstrated for NV1 \\
		Apply the method to a real object & SRIO instantiation; stored evidence & demonstrated \\
		Distinguish applicable, blocked, and open tests & four-outcome status model; Tables \ref{tab:cm_applicability} and \ref{tab:req_verification} & demonstrated \\
		Support integration into future releases & lifecycle and governance model; reusable schema & specified by design \\
		\bottomrule
	\end{tabularx}
\end{table}

The empirical validation covered the central methodology steps under NV1. The requirement derivation, the asset mapping, the applicability assessment, the test selection, the execution, and the evidence-based assessment were all exercised on a real device. The tests that require passive observation or inline manipulation remain extensions of the validated core process. The methodology is therefore validated in part, within a defined boundary at the NV2 and NV3 variants and at the requirements whose evidence is still open. A claim of full validation is not made. The validation concerns the methodology rather than the product, and its boundary is the NV1 position together with the requirements whose evidence remains open; these limits are examined further in Section \ref{subsec:threats_to_validity}.

\subsection{Reliability and Threats to Validity}
\label{subsec:threats_to_validity}

The results are subject to threats to construct validity, internal validity, external validity and reliability.

\textbf{Construct validity.} The fourteen requirements are an interpretation of machine-level obligations and were derived by a single author. They are mapped through the draft standard prEN 50742 to IEC 62443-4-2, so the interpretation may change if the final standard text differs. The mapping is therefore itself a threat to validity. Several practical observations are also indirect: the Nikto Actuator finding remains unconfirmed, a bad-request response does not distinguish input validation from authorisation, and the SSLscan result establishes the absence of a TLS endpoint but not a complete confidentiality assessment.

\textbf{Internal validity.} The fuzzing-induced HTTP disruption was not reproduced, and its causation cannot be attributed to a specific input. The layer-two flooding observation lacks a complete terminal capture. In addition, hping3 and macof target the switch and the shared segment as well as the device, so the observed fault indication cannot be attributed to a device-side security response.

\textbf{Safety-priority constraint.} Chapter 3 requires that testing must not impair the safety function. Under NV1, however, the cyclic safety channel was not observable. The constraint was therefore maintained procedurally through the isolated laboratory and the exclusion of production systems rather than demonstrated by measurement.

\textbf{Connection-hold result.} No stored evidence establishes that the HTTP connection slots were occupied during the hold procedure. Its negative result is consequently inconclusive and does not demonstrate resilience.

\textbf{External validity.} The evaluation covers one device, one firmware revision, one flat laboratory topology without VLAN separation, and only the NV1 network position. The results do not generalise to other gateways or to the cyclic safety path.

\textbf{Reliability.} The work was performed by a single tester and depends on specific tool versions and scan timing. The Nmap timing sensitivity is itself evidence that execution conditions affect observations. The stored tool outputs, commands in Appendix \ref{app:tool_commands}, and defined acceptance criteria make the raw evidence reproducible, although the interpretation is not fully reproducible where causation or setup state is uncertain.

These limitations bound the claims of Section \ref{subsec:objective_validation_sub} to a partial validation of the methodology.

% TODO-OTHERFILE: verify that the abstract, the German Kurzfassung and the summary chapter state the same counts as Chapter 4 after these changes: 14 derived requirements; 9 assets; 10 company tests; 5 applicable, 2 not applicable, 2 executed with negative result and 1 blocked; 15 executed test activities; 2 FAIL verdicts (RQ-003 and RQ-012); 4 NOT ASSESSED verdicts (RQ-002, RQ-005, RQ-009 and RQ-013).



